\chapter*{기호와 용어}
\addcontentsline{toc}{chapter}{기호와 용어}

\renewcommand{\arraystretch}{1.2}
\begin{longtable}{>{\raggedright\arraybackslash}p{0.25\textwidth}>{\raggedright\arraybackslash}p{0.65\textwidth}}
\toprule
\textbf{기호·용어} & \textbf{뜻} \\
\midrule
\endfirsthead
\toprule
\textbf{기호·용어} & \textbf{뜻} \\
\midrule
\endhead
\bottomrule
\endfoot

$L/K$ & 체 $K$를 부분체로 포함하는 체 $L$의 확장. \\
$[L:K]$ & $K$-벡터공간 $L$의 차원, 즉 확장의 차수. \\
중간체 & $K\subseteq E\subseteq L$을 만족하는 부분체 $E$. \\
$\charac K$ & 체 $K$의 표수. \\
소체 & 체 안의 가장 작은 부분체. 표수 $0$이면 $\Q$, 표수 $p$이면 $\F_p$와 동형이다. \\
$\Frob_K$ & 표수 $p$인 체에서 $a\mapsto a^p$로 주어지는 프로베니우스 준동형. \\
대수적 원소 & 영이 아닌 $K$-계수 다항식의 근이 되는 원소. \\
초월원소 & $K$ 위에서 어떤 영이 아닌 다항식 관계도 만족하지 않는 원소. \\
$m_{\alpha,K}$ & 원소 $\alpha$의 체 $K$ 위 최소다항식. \\
$K[\alpha]$ & $K$와 $\alpha$가 생성하는 부분환. \\
$K(\alpha)$ & $K$와 $\alpha$가 생성하는 부분체. \\
단순확장 & 어떤 $\alpha$에 대해 $L=K(\alpha)$인 체확장. \\
원시원 & 유한확장 $L/K$에서 $L=K(\theta)$를 만족하는 생성원 $\theta$. \\
$K(\alpha_1,\dots,\alpha_r)$ & $K$와 원소 $\alpha_1,\dots,\alpha_r$가 생성하는 가장 작은 부분체. \\
탑 법칙 & $K\subseteq L\subseteq M$에서 $[M:K]=[M:L][L:K]$. \\
$\overline K$ & 체 $K$의 대수적 폐포. \\
분해체 & 주어진 다항식이 완전히 분해되고 그 모든 근으로 생성되는 최소 확장체. \\
대수적으로 닫힌 체 & 모든 비상수 다항식이 체 안에 근을 갖는 체. \\
분리다항식 & 대수적 폐포에서 중근을 갖지 않는 다항식. \\
분리원소·분리확장 & 최소다항식이 분리인 원소와 모든 원소가 분리적인 대수적 확장. \\
$\Hom_K(L,\overline K)$ & $K$를 고정하는 체매장 $L\to\overline K$의 집합. \\
$[L:K]_{\mathrm s}$ & $\#\Hom_K(L,\overline K)$로 정의되는 분리차수. \\
완전체 & 모든 대수적 확장이 분리확장인 체. 표수 $p$에서는 프로베니우스가 전사인 것과 동치이다. \\
순수비분리 원소·확장 & 충분한 $p$-거듭제곱이 기준체에 들어가는 원소와 그러한 원소들로 이루어진 대수적 확장. \\
$K_{\mathrm s}$ & 대수적 확장 $L/K$ 안에서 $K$ 위 분리적인 모든 원소가 이루는 분리폐포. \\
$[L:K]_{\mathrm i}$ & 유한확장에서 $[L:K]/[L:K]_{\mathrm s}$로 정의되는 비분리차수. \\
$\overline\Q_{\mathrm{alg}}$ & $\C$ 안에서 $\Q$ 위 대수적인 원소들의 체. \\

\end{longtable}
